\begin{lemma}
Let $T$ be finite, $J\subseteq T$, and $0\le p\le1$. Put
$\rho=(p\delta_1+(1-p)\delta_0)\otimes(\operatorname{Leb}|_{[0,1]})$,
$\mu=\rho^T$, $\tau=\rho^{T\setminus J}$, and $\kappa=\rho^J$,
all on the natural finite-dimensional Borel spaces. Define
$F(t,c)_i=c_i$ for $i\in J$ and $F(t,c)_i=t_i$ otherwise.
The map $(t,c)\mapsto F(t,c)$ and the section $t\mapsto F(t,0)$ are
measurable. The measures $\mu,\tau,\kappa$ each have total mass one,
and $\tau$-almost every $t$ has every priority coordinate in $[0,1]$.
For every measurable $E\subseteq(\mathbb R^2)^T$, the function
$t\mapsto\kappa\{c:F(t,c)\in E\}$ is measurable and
\[
 \mu(E)=\int \kappa\{c:F(t,c)\in E\}\,d\tau(t).
\]
The integral is the nonnegative extended integral.
\end{lemma}
\begin{proof}
For a fixed index $i$, evaluation of $F$ is either a current-coordinate
projection or a context-coordinate projection; the choice depends only on
$i\in J$. Both are measurable, so the finite product map $F$ is
measurable. Substituting the constant zero current tuple proves the section
assertion. Restriction to the two complementary index subsets is a
measurable inverse to $F$.

The activation factor has total mass $p+(1-p)=1$ and the restricted
Lebesgue factor has mass one. The finite product rectangle formula then
gives total mass one for all three measures, even when one index subset
is empty. Each priority marginal is supported on $[0,1]$, since it is the
restricted Lebesgue factor and all other factors have mass one. The
complement of simultaneous priority support is a finite union of these
coordinate null sets, hence is null.

We verify $F_*(\tau\otimes\kappa)=\mu$. For a Borel rectangle
$E=\prod_{i\in T}E_i$, its inverse image under $F$ is
$(\prod_{i\notin J}E_i)\times(\prod_{i\in J}E_i)$. The product
rectangle formula assigns this inverse image mass
\[
 \left(\prod_{i\notin J}\rho(E_i)\right)
 \left(\prod_{i\in J}\rho(E_i)\right)=\prod_{i\in T}\rho(E_i)=\mu(E).
\]
These rectangles form a generating intersection-stable class containing
the whole space. Both measures are finite with the same total mass; the
finite-measure uniqueness theorem (equivalently the pi-lambda theorem)
therefore extends the equality to the product Borel sigma algebra.
This is precisely the coordinate-splitting measure-preserving identity.

For a measurable $E$, the indicator $1_E(F(t,c))$ is nonnegative and
jointly measurable. Tonelli's theorem for the two finite measures makes
its current-coordinate integral a measurable function of $t$, and gives
\[
 \mu(E)=\int 1_E(F(t,c))\,d(\tau\otimes\kappa)(t,c)
       =\int\left(\int 1_E(F(t,c))\,d\kappa(c)\right)d\tau(t).
\]
The inner integral is the measure of the displayed fiber. This proves
both conclusions, without regular conditional probabilities or conditioning
on any event of a prescribed continuous coordinate value.
\end{proof}
